% CV Harvard — Axel Pfingsten Arpe
% Compilar: pdflatex CV_Axel_Pfingsten_Arpe.tex
\documentclass[11pt,a4paper]{article}

\usepackage[utf8]{inputenc}
\usepackage[T1]{fontenc}
\usepackage[spanish,es-nodecimaldot]{babel}
\usepackage{mathptmx} % Times (estilo Harvard clásico)
\usepackage[margin=1.6cm]{geometry}
\usepackage{enumitem}
\usepackage{titlesec}
\usepackage{hyperref}
\usepackage{xcolor}
\usepackage{setspace}

\definecolor{linkblue}{RGB}{0,82,147}
\hypersetup{
  colorlinks=true,
  urlcolor=linkblue,
  linkcolor=linkblue,
  pdfauthor={Axel Pfingsten Arpe},
  pdftitle={Curriculum Vitae — Axel Pfingsten Arpe},
  pdfsubject={Enfermero Profesional — Intensivo / continuidad}
}

\pagestyle{empty}
\setlength{\parindent}{0pt}
\setlength{\parskip}{0pt}
\setlist[itemize]{leftmargin=1.2em, itemsep=1.5pt, parsep=0pt, topsep=2pt}

% Encabezados de sección estilo Harvard
\titleformat{\section}
  {\large\bfseries\uppercase}
  {}{0pt}{}[\titlerule]
\titlespacing*{\section}{0pt}{10pt}{6pt}

\newcommand{\cventry}[3]{%
  \textbf{#1} \hfill #2\\
  \textit{#3}\\[2pt]
}

% Enlace a Google Maps (el texto visible se mantiene para ATS e impresión)
\newcommand{\maps}[2]{%
  \href{https://www.google.com/maps/search/?api=1\&query=#1}{#2}%
}

\begin{document}
\setstretch{1.05}

%========== ENCABEZADO ==========
\begin{center}
  {\LARGE\bfseries Axel Pfingsten Arpe}\\[4pt]
  {\large Enfermero Profesional --- Cuidados intensivos / continuidad del cuidado}\\[6pt]
  San Joaquín, Región Metropolitana, Chile\\[2pt]
  +56\,9\,4203\,5552 \quad$\vert$\quad
  \href{mailto:axel.pfingsten@gmail.com}{axel.pfingsten@gmail.com}\\[2pt]
  \href{https://www.linkedin.com/in/axel-pfingsten-arpe-43781a209/}{Axel Pfingsten}\\[2pt]
  \textbf{Disponibilidad inmediata}
\end{center}

\vspace{2pt}

%========== PERFIL ==========
\section{Perfil profesional}
Enfermero titulado con experiencia en SAC/SAR, procedimientos e IAAS, y en continuidad de pacientes de mayor complejidad en APS. Cursos IAAS (80~h) y RCP/DEA vigentes. SIS al día; no acredito experiencia en UCI. Disponibilidad inmediata para turnos en Providencia.

%========== COMPETENCIAS ==========
\section{Competencias clave}
\textbf{Clínica:} SAC/SAR, continuidad, mayor complejidad.\\[2pt]
\textbf{Procedimientos:} curaciones, medicación, Foley, IAAS, RCP/DEA.\\[2pt]
\textbf{Gestión:} RAYEN, calidad, trabajo en equipo.\\[2pt]


%========== EXPERIENCIA ==========
\section{Experiencia profesional}

\cventry{Enfermero de Sector}
{Nov 2025 -- Ene 2026}
{\maps{CESFAM+La+Florida+Av.+La+Florida+6015,+La+Florida,+Santiago}{CESFAM La Florida — La Florida, RM}}
\begin{itemize}
  \item Atendí en box de APS, coordinando seguimiento y funcionamiento del sector.
  \item Coordiné al equipo en curaciones y pie diabético, asegurando calidad del servicio.
  \item Registré evolución clínica, asegurando continuidad y calidad de la atención.
\end{itemize}
\vspace{6pt}

\cventry{Enfermero de Sector}
{Ago 2024 -- Nov 2025}
{\maps{CESFAM+Bellavista+Pudeto+7100,+La+Florida,+Santiago}{CESFAM Bellavista de La Florida — La Florida, RM}}
\begin{itemize}
  \item Coordiné ingreso, seguimiento y evaluación de NANEAS, ECICEP y pacientes postrados.
  \item Organicé atención de sector con enfoque familiar, coordinando el plan de cuidado del equipo.
  \item Realicé controles de niño sano (2 meses a 9 años), evaluaciones EEDP y TEPSI, y clínicas de lactancia.
  \item Apliqué evaluación de pie diabético, insulinoterapia y ajuste de tratamiento en hipertensión arterial (HTA).
  \item Eduqué en autocuidado y prevención, alineado a promoción de salud en APS.
\end{itemize}
\vspace{6pt}

\cventry{Enfermero Programa de Pacientes Postrados}
{May 2024 -- Ago 2024}
{\maps{CESFAM+N5+Union+Latinoamericana+98,+Santiago+Centro}{CESFAM N\textdegree5 — Santiago Centro, RM}}
\begin{itemize}
  \item Evalué integralmente a pacientes postrados en domicilio, coordinando alertas con el equipo.
  \item Ejecuté curaciones, medicación e inyectables en domicilio, educando a cuidadores.
  \item Realicé toma de muestras y Foley en terreno, con registro clínico actualizado.
  \item Coordiné continuidad del cuidado mediante informes claros al equipo interdisciplinario.
\end{itemize}
\vspace{6pt}

\cventry{Enfermero de Sector}
{Ene 2022 -- Dic 2023}
{\maps{CESFAM+Dr.+Fernando+Maffioletti+Av.+Central+301,+La+Florida,+Santiago}{CESFAM Dr.\ Fernando Maffioletti — La Florida, RM}}
\begin{itemize}
  \item Desempeñé funciones en SAC, áreas respiratorias y COVID-19, con foco en continuidad clínica.
  \item Ejecuté procedimientos (curaciones, medicación y Foley) aplicando protocolos de IAAS.
  \item Realicé toma de muestras con identificación/registro, más test de antígeno y PCR.
  \item Atendí en box APS: controles de niño sano, EEDP, TEPSI y PSCV.
  \item Coordiné seguimiento de pacientes crónicos y postrados, asegurando continuidad del cuidado.
\end{itemize}
\vspace{6pt}

\cventry{Enfermero Área Box}
{Nov 2021 -- Dic 2021}
{\maps{CESFAM+Santa+Amalia+202,+La+Florida,+Santiago}{CESFAM Santa Amalia — La Florida, RM}}
\begin{itemize}
  \item Atendí en box PSCV y controles de niño sano (2 meses a 9 años).
  \item Apliqué EMPA/EMPAM y eduqué en prevención y autocuidado.
  \item Notifiqué y di seguimiento a casos COVID-19 en área respiratoria.
\end{itemize}
\vspace{6pt}

\cventry{Enfermero de Apoyo APS}
{Abr 2021 -- Oct 2021}
{\maps{CESFAM+Los+Quillayes+Julio+Cesar+10905,+La+Florida,+Santiago}{CESFAM Los Quillayes — La Florida, RM}}
\begin{itemize}
  \item Apoyé la unidad de procedimientos y SAR, asegurando recursos y esterilidad del box.
  \item Ejecuté curaciones, medicación y dispositivos invasivos con foco en bioseguridad.
  \item Apliqué protocolos de IAAS y calidad en cada procedimiento clínico.
  \item Participé en campañas de vacunación COVID-19 y VPH, fortaleciendo cobertura del PNI.
  \item Realicé atención en box: controles preventivos y control sano infantil.
\end{itemize}

%========== FORMACIÓN ==========
\section{Formación académica}

\cventry{Título profesional de Enfermería}
{Enero 2021}
{Universidad Central de Chile}

\vspace{4pt}
\textbf{Diplomado y cursos de especialización} — AGCS Asesorías y Capacitación\\[2pt]
\begin{itemize}
  \item Diplomado de Calidad y Gestión en Salud (320 horas) — Julio 2026
  \item Curso Prevención y Control de IAAS (80 horas) — Mayo 2026
  \item Curso Reanimación Cardiopulmonar (RCP) y DEA (30 horas) — Mayo 2026
  \item Curso Programa Nacional de Inmunizaciones — PNI (60 horas) — Junio 2026
  \item Curso Abordaje Integral del Trastorno del Espectro Autista — TEA (80 horas) — Junio 2026
\end{itemize}

%========== REFERENCIAS ==========
\section{Referencias}
\begin{itemize}
  \item \textbf{Giannina Cavallo Barría} — Directora, CESFAM Maffioletti\\
        \href{mailto:Gcavallob@gmail.com}{Gcavallob@gmail.com} \quad$\vert$\quad +56\,9\,8200\,8584
  \item \textbf{Claudio Andrés Silva Orellana} — Director, CESFAM Los Quillayes\\
        \href{mailto:caso.klgo@gmail.com}{caso.klgo@gmail.com} \quad$\vert$\quad +56\,9\,9079\,9877
  \item \textbf{Silvia Acuña Vielma} — Encargada de programa infantil, CESFAM Maffioletti\\
        \href{mailto:silvia.avielma@gmail.com}{silvia.avielma@gmail.com} \quad$\vert$\quad +56\,9\,5015\,5455
  \item \textbf{María Dolores Barja Aguilera} — Jefa de sector, CESFAM Bellavista de La Florida\\
        \href{mailto:valle.bv@comudef.cl}{valle.bv@comudef.cl} \quad$\vert$\quad +56\,9\,8549\,1392
  \item \textbf{Carla Esquivel Rozas} — Subdirectora, CESFAM La Florida\\
        \href{mailto:subdireccioncslf@comudef.cl}{subdireccioncslf@comudef.cl} \quad$\vert$\quad +56\,9\,9162\,7012
\end{itemize}

\end{document}
